\problemname{Compressing Commands}

\illustration{0.3}{paths.jpg}{
    Only true terminal-dwellers use green on black colours.
    CC BY-NC 2.0 by Kjetil Korslien on \href{https://www.flickr.com/photos/kjetikor/8483021009/}{Flickr}
}

% optionally define variables/limits for this problem
\newcommand{\maxn}{10^5}
\newcommand{\maxs}{10^6}

Unlike your friends, you live in a terminal. You are cool. Your terminal is
everything to you: you have optimized the font, colour scheme, keyboard
shortcuts, and what not.

Once thing still annoys you though: all these commands you're typing involve so
many file paths and are so long! Then it occurs to you: all this time you have
been working from the \emph{root directory} of the file system, but in fact you
can \emph{change directory} to anywhere you like! This should simplify your life
a lot!

This way, if your \emph{working directory} is \texttt{/a/b}, you can refer to
the absolute file path \texttt{/a/b/x} using simply \texttt{x}. To \emph{go up} a
level, you can use \texttt{..}, so that you can refer to \texttt{/a/y/z}
as \texttt{../y/z}, and to \texttt{/some/other/directory} as \texttt{../../some/other/directory}.

You being you, of course you overdo this and will now use relative
paths \emph{everywhere}!

Given the $n$ absolute file paths in the command you want to run, find the
working directory that minimizes the total number of relative path components.
For example, \texttt{a/a/b/c}, and \texttt{../../a/b} both contain $4$ path components.
Note that you can only change the working directory to a directory and not to a file path.
Filenames will never coincide with directory names in the same directory.

In the first sample it is best to set the working directory
to \texttt{/home/jury/compressingcommands}, leading to $6$
components: \texttt{secret}, \texttt{solutions},
and \texttt{../../hackerman/answers}.

In the second sample it is best to set the working directory to \texttt{/a},
leading to $19$ path components: \texttt{b/a/a/b}, \texttt{a/a}, \texttt{../b}, \texttt{a/b}, \texttt{b/a/a/a}, \texttt{../c}, and \texttt{b/a/b}.

\begin{Input}
    The input consists of:
    \begin{itemize}
        \item One line with an integer $n$ ($1\leq n\leq \maxn$), the number of
    absolute file paths.
        \item $n$ lines, each with a string $s$, an absolute file path. Each path
    contains only lowercase English letters (\texttt{a-z}) and slashes
    (`\texttt{/}'), starts with `\texttt{/}', does not end with `\texttt{/}', and does not
    contain consecutive slashes (``\texttt{//}'').
    \end{itemize}
    The total number of characters in the $n$ strings is at most $\maxs$.
\end{Input}

\begin{Output}
    Output the minimal number of relative path components that can be achieved by
    changing to a different working directory.
\end{Output}
